%% toc: P1 4 4.1,4.2,4.3,4.4
\section{Proofs for the numerical coding of strings}
\label{sec:annex-p01-w02-p01-012-string-coding}

Throughout, \(e:A\to\N\) is an injection, the stage codes \(e_n\) are those
of \cref{def:p01-w02-p01-012-coding-maps}, and \(t_s=s(s+1)/2\).

\begin{annexitem}{annex:p01-w02-p01-012-stage-injective}{thm:p01-w02-p01-012-stage-injective}{Every stage code is injective}
\begin{proof}
\emph{Step 0: the property of the pairing map that is used.} If
\(\operatorname{pair}_{\N}(a,b)=\operatorname{pair}_{\N}(a',b')\), then
\(a=a'\) and \(b=b'\). This is the injectivity half of the bijectivity
argument of \cref{sec:strings-01}; it is written out here in the form the
induction uses. Each \(t_s\) is a natural number, since one of \(s\) and
\(s+1\) is even, and \(t_{s+1}-t_s=\bigl((s+1)(s+2)-s(s+1)\bigr)/2=s+1\), so
\(t_0<t_1<t_2<\cdots\). Let \(z=\operatorname{pair}_{\N}(a,b)\) and
\(s=a+b\), so that \(z=t_s+b\). Then \(t_s\leq z\) because \(b\geq0\), and
\(z=t_s+b\leq t_s+s<t_s+s+1=t_{s+1}\) because \(b\leq a+b=s\). Likewise,
with \(s'=a'+b'\), \(t_{s'}\leq z<t_{s'+1}\). Suppose \(s<s'\). Then
\(s+1\leq s'\), and since the \(t_k\) increase, \(t_{s+1}\leq t_{s'}\); this
gives \(z<t_{s+1}\leq t_{s'}\leq z\), which is impossible. By symmetry
\(s'<s\) is impossible as well, so \(s=s'\). Then \(b=z-t_s=z-t_{s'}=b'\) and
\(a=s-b=s'-b'=a'\).

\emph{Base case, \(n=0\).} A sequence of length \(0\) is a function with
domain \(\{0,\ldots,-1\}=\emptyset\). There is exactly one function with
empty domain, so \(A^0=\{\epsilon\}\). A map on a set with one element is
injective: if \(e_0(w)=e_0(w')\) with \(w,w'\in A^0\), then
\(w=\epsilon=w'\).

\emph{Inductive step.} Assume that \(e_n\) is injective, and let
\(w,w'\in A^{n+1}\) satisfy \(e_{n+1}(w)=e_{n+1}(w')\). By the definition of
\(e_{n+1}\),
\[
  \operatorname{pair}_{\N}\bigl(e_n(w\mathbin{\upharpoonright}n),e(w(n))\bigr)
  =\operatorname{pair}_{\N}\bigl(e_n(w'\mathbin{\upharpoonright}n),e(w'(n))\bigr).
\]
Step 0 applied with \(a=e_n(w\mathbin{\upharpoonright}n)\),
\(b=e(w(n))\), \(a'=e_n(w'\mathbin{\upharpoonright}n)\) and
\(b'=e(w'(n))\) gives
\[
  e_n(w\mathbin{\upharpoonright}n)=e_n(w'\mathbin{\upharpoonright}n)
  \qquad\text{and}\qquad
  e(w(n))=e(w'(n)).
\]
Both restrictions lie in \(A^n\), so the induction hypothesis gives
\(w\mathbin{\upharpoonright}n=w'\mathbin{\upharpoonright}n\); that is,
\(w(i)=w'(i)\) for every \(i<n\). The letters \(w(n)\) and \(w'(n)\) lie in
\(A\), so the injectivity of \(e\) gives \(w(n)=w'(n)\). Hence \(w\) and
\(w'\) have the same domain \(\{0,\ldots,n\}\), which is
\(\{0,\ldots,n-1\}\cup\{n\}\), and agree at every argument in it. Two
functions with the same domain and the same values are equal, so \(w=w'\).
Thus \(e_{n+1}\) is injective.

By induction on \(n\), every \(e_n\) is injective.
\end{proof}
\end{annexitem}

\begin{annexitem}{annex:p01-w02-p01-012-tagged-injective}{thm:p01-w02-p01-012-tagged-injective}{The length-tagged code is an injection}
\begin{proof}
\emph{Part 1.} Let \(w\in A^m\cap A^n\). Then \(w\) is a function with domain
\(\{0,\ldots,m-1\}\) and also a function with domain \(\{0,\ldots,n-1\}\). A
function has exactly one domain, so these two sets are equal. The first has
exactly \(m\) elements and the second exactly \(n\), so \(m=n\). Hence for
\(m\neq n\) no such \(w\) exists, and \(A^m\cap A^n=\emptyset\). Each
\(w\in A^*=\bigcup_{n\in\N}A^n\) lies in some \(A^n\), and by what was just
shown in no other; this \(n\) is \(\operatorname{len}(w)\). The rule
\(w\mapsto\operatorname{pair}_{\N}(n,e_n(w))\) for \(w\in A^n\) therefore
assigns exactly one value to each \(w\in A^*\), so it is a function
\(A^*\to\N\).

\emph{Part 2.} Let \(w\in A^n\) and \(w'\in A^{n'}\), and suppose
\(\operatorname{pair}_{\N}(n,e_n(w))=\operatorname{pair}_{\N}(n',e_{n'}(w'))\).
Step 0 of the proof of \cref{thm:p01-w02-p01-012-stage-injective}, applied
with \(a=n\), \(b=e_n(w)\), \(a'=n'\) and \(b'=e_{n'}(w')\), gives \(n=n'\)
and \(e_n(w)=e_{n'}(w')=e_n(w')\). Now \(w\) and \(w'\) both lie in \(A^n\),
and \(e_n\) is injective by \cref{thm:p01-w02-p01-012-stage-injective}, so
\(w=w'\). Hence equal tagged codes come from equal strings, and the
length-tagged code is an injection \(A^*\to\N\).
\end{proof}
\end{annexitem}

\begin{annexitem}{annex:p01-w02-p01-012-countable}{thm:p01-w02-p01-012-countable}{Countability of finite strings}
\begin{proof}
Let \(A\) be countable. By the definition of countable set in
\cref{sec:set-prerequisites-05}, there is an injection \(e:A\to\N\); fix one.
Define the stage codes \(e_n\) from this \(e\) by
\cref{def:p01-w02-p01-012-coding-maps}. By
\cref{thm:p01-w02-p01-012-stage-injective} every \(e_n\) is injective, and by
\cref{thm:p01-w02-p01-012-tagged-injective} the length-tagged code is an
injection \(A^*\to\N\). So \(A^*\) is countable, again by the definition of
countable set.
\end{proof}
\end{annexitem}

\begin{annexitem}{annex:p01-w02-p01-012-decoding-correct}{thm:p01-w02-p01-012-decoding-correct}{Decoding recovers the string}
\begin{proof}
\emph{The inverse of the pairing map.} If \(z=\operatorname{pair}_{\N}(a,b)\),
then \(\operatorname{unpair}(z)=(a,b)\). Indeed, Step 0 of the proof of
\cref{thm:p01-w02-p01-012-stage-injective} shows that
\(t_{a+b}\leq z<t_{a+b+1}\), and the index \(s\) with \(t_s\leq z<t_{s+1}\)
is unique because the \(t_k\) increase (two such indices \(s<s'\) would give
\(z<t_{s+1}\leq t_{s'}\leq z\)). So \(s=a+b\), the second inverse value is
\(z-t_{a+b}=b\), and the first is \(s-b=a\).

\emph{Induction on \(n\): \(d_n(e_n(w))\) is defined and equals \(w\) for
every \(w\in A^n\).} For \(n=0\), the only element of \(A^0\) is \(\epsilon\)
(base case of \cref{thm:p01-w02-p01-012-stage-injective}), and
\(d_0(e_0(\epsilon))=d_0(0)=\epsilon\).

Assume the claim for \(n\), and let \(w\in A^{n+1}\) and
\[
  y=e_{n+1}(w)=\operatorname{pair}_{\N}\bigl(e_n(w\mathbin{\upharpoonright}n),
  e(w(n))\bigr).
\]
By the previous paragraph,
\(\operatorname{unpair}(y)=\bigl(e_n(w\mathbin{\upharpoonright}n),
e(w(n))\bigr)\). The number \(b=e(w(n))\) is a value of \(e\), and the unique
letter \(x\) with \(e(x)=b\) is \(x=w(n)\), unique because \(e\) is
injective. Since \(w\mathbin{\upharpoonright}n\in A^n\), the induction
hypothesis gives that
\(d_n\bigl(e_n(w\mathbin{\upharpoonright}n)\bigr)\) is defined and equals
\(w\mathbin{\upharpoonright}n\). Hence \(d_{n+1}(y)\) is defined and
\[
  d_{n+1}(y)=(w\mathbin{\upharpoonright}n)\langle w(n)\rangle .
\]
By the definition of concatenation in \cref{sec:strings-01}, the right side
is the sequence of length \(n+1\) whose value at \(i<n\) is
\((w\mathbin{\upharpoonright}n)(i)=w(i)\) and whose value at \(n\) is
\(\langle w(n)\rangle(n-n)=w(n)\). It has the domain and the values of \(w\),
so it equals \(w\). This completes the induction.

\emph{The procedure.} Let \(w\in A^n\) and
\(z=\operatorname{pair}_{\N}(n,e_n(w))\). Step 1 computes
\(\operatorname{unpair}(z)=(n,e_n(w))\) by the first paragraph, and step 2
computes \(d_n(e_n(w))\), which by the induction is defined and equals \(w\).
So the procedure outputs \(w\).
\end{proof}
\end{annexitem}
